\documentclass[book.tex]{subfiles}
% Merge with case_study_molecularProperties
\begin{document}

\chapter{Case Studies in Material Science}

\label{ch:ml_ms}
\noindent\rule{11cm}{0.4pt} \\
\textbf{Prerequisites:} \autoref{chap:graphconvs}, \autoref{chap:molecular_ml}   \\
\textbf{Difficulty Level:} ** \\
\noindent\rule{11cm}{0.4pt}
\vspace{5mm}

In this chapter, we walk through two case studies building deep learning models for predicting collective properties of materials. The first case study focuses on predicting melting and boiling points of liquids, which is relevant in several applications. For example, in electrochemical storage and conversion devices, a design condition involving both melting and boiling points for the electrolyte composition is the existence in liquid phase through the desirable temperature range of operation. The second case study focuses on predicting solubility of salts and additives in liquids, which is relevant in every multi-component liquid-phase material for a variety of application areas.

\section{Melting and Boiling Points of Liquids}

An example application where rapid screening is necessary is that of electrochemical systems such as batteries. Liquid electrolytes in batteries generally consist of a lithium salt dissolved in a solvent. These electrolytes allow for the transport of Li ions which is the key performance variable determining Li-ion batteries. At present, a chemical class known as carbonates are the commonly used solvents, but they suffer from sensitivity to temperature and are volatile and flammable \cite{Yamada2019AdvancesElectrolytes}.  There is an urgent need to find electrolyte solvents that are non-flammable and have a wider temperature range of operation.

This presents a need to identify new solvents that can carry novel electrochemical systems forward. Screening of these solvents to hone in on the most promising candidates is desired to speed up the discovery process. In particular, urgency towards addressing the current climate crisis means that fast screening could help meet the immediate need for improved green technologies. Choice of solvent employed in the electrolyte can have significant impact on overall cell performance, and an extensive list of criteria for the ideal candidate emerges \cite{Ue2014}. One of these criteria is the size of the operating temperature window of the cell, for which the solvent plays a large role. It is necessary that the electrolyte remains a liquid during operation to provide sufficient conductivity.  Close to its melting point, the liquid becomes viscous and the conductivity drops precipitously.  Thus, this places restrictions on the allowed melting points and boiling points of the electrolyte solvent.  Given the richness of organic chemistry (millions of possible compounds), experimental testing of each compound is nearly impossible. Additionally, as new solvents are synthesized, their properties such as these critical points may not be immediately available. 

Prediction of melting and boiling points using statistical mechanics is a challenging procedure and would represent a slow screening process. On the other hand, utilizing the recent advances within machine learning techniques allows for the leapfrogging from the molecular structure to the melting and boiling points. In particular, innovations in molecular featurization techniques have opened the door to encoding chemical structure along with atomic features \cite{Wu2018MoleculeNet:Learning}. In this case study, we show and contrast the use of extended connectivity fingerprints (ECFP), graph convolutional neural network models (GCNN), and Weave models to simultaneously predict the melting and boiling points of $\approx$3900 substances from the molecular structure of each. The dataset used for this case study was collected with no restrictions on their material class thereby giving a general set of materials. Data was collected from the Pubchem which is the world’s largest open chemistry database \cite{10.1093/nar/gky1033} by parsing over 500,000 compound identifiers to identify materials for which both boiling and melting points were reported.


We show that ECFP was unable to provide satisfactory predictions in this context, but both GCNN and Weave demonstrated promising predictive capabilities, with the Weave model providing the best performance when testing. 

% \subsection{Data}

%  This database comprises of molecules uploaded by authors after publication of their articles or patents, and is maintained by the National Institutes of Health. Over 97 million compounds are included with over 237 million substances stemming from over 30 million papers and over 3 million patents. For each of these materials melting and boiling points are provided for many (but not all) in addition to the molecular structure and a unique SMILES identifier. For this case study, only materials containing reports for both Melting and Boiling points were used. Applying this condition, 3909 materials were identified after parsing over 500,000 compound identifiers. Further details regarding the method of parsing are provided in the Methods section.

% It was found when training that including all data points yielded poor root mean squared errors (RMSE). When observing the range of temperatures of the samples (see Fig. \ref{fig:data_hist}), it becomes apparent that very few data points have boiling points that lie above 700 K (3.8\%), so these were regarded as outliers and removed from the following analysis. More discussion on this is included in the Weave evaluation section.


%\begin{marginfigure}[-10cm]
%\begin{center}
% \fbox{\rule{0pt}{2in} %\rule{0.9\linewidth}{0pt}
%   \includegraphics[width=0.8\linewidth]{figures/hist_data_final.png}}
%\end{center}
%   \caption{Histogram of critical temperatures of all ($\approx$3900) data points. A small fraction (3.8\%) of the total data with critical temperatures higher than 700K was discarded.}
%\label{fig:data_hist}
%\end{marginfigure}


%------------------------------------------------------------------------

The overall workflow for this case study can be broken down into two stages: data collection and model optimization.

%\begin{marginfigure}
%\begin{center}
%\fbox{\rule{0pt}{2in} %\rule{0.9\linewidth}{0pt}
%   \includegraphics[width=0.8\linewidth]{figures/work_flow.png}}
%\end{center}
%   \caption{Overall workflow of the case study with the paths of data collection and featurization converging at the model selection and training}
%\label{fig:workflow}
%\end{marginfigure}

% \textit{\textbf{Data Collection}}

% Under the Experimental Properties section of Pubchem, one can obtain the boiling and melting point information of any given compound reported by many different literature sources. To store molecular data for such a large number of compounds and avoid manual data collection, two different python libraries were used to access compound IDs within PubChem. This was done using two libraries (PubChemPy \cite{swain2014pubchempy} and Beautiful Soup \cite{richardson2007beautiful}) which are used to get the input molecular structures IDs and experimental properties of the compounds respectively.

% The simplified molecular-input line-entry system identity (SMILES ID) uses short ASCII strings to describe the structure of a compound to be fed into different featurizers as input data for the model. PubChempy allows chemical searches by their compound IDs (CIDs) and returns different properties of the compound including the SMILES ID. It is fast in terms of extracting the features as it uses chemical toolkits by a Univeral Resource Locator (URL) based application program interface (API) called the PUG REST (Power User Gateway Representational State Transfer) web service. 

% At present, PubChemPy is not capable of extracting second order data stored within sub headings of a particular compound like Boiling point within Experimental Properties. Since all information for a PUG-REST request is encoded into its URL, we parsed individual eXtensible Markup Language (XML) files of these URLs using another python library called the Beautiful Soup. This library is primarily employed in programmatic extraction of data from HTML and XML files using python language. By pointing to the required sub heading in the PUG REST URL, we could easily obtain strings containing MP or BP temperatures.


\subsection{Models} 
The Deepchem package \cite{ramsundar2019deep} was used in this case study to employ all the three featurization techniques discussed. These three techniques all share the property of working to include structural features into the models \cite{Wu2018MoleculeNet:Learning}. This is particularly important in the context of melting and boiling points as they are governed by interatomic interactions. For each of the featurization techniques and models studied, 5-fold cross validation was used to avoid over-optimizing towards a given subset of the data and to optimize hyperparameters. 20 \% of the total data was set aside for testing, and the remaining data used an 80-20 split for training and validation. As we are interested in predicting both melting and boiling points, we used multitask models to simultaneously train for both quantities. Hypothetically, both melting and boiling points should be influenced by similar atomic features, particularly those related to bonding strength. %Here we outline each of the three approaches in greater detail, and highlight their distinguishing features.

%Extended-Connectivity Fingers (ECFP): This featurizer contains topological characteristics of the molecule and is commonly applied to tasks such as similarity searching and activity prediction. During featurization, the molecule is divided into submodules which are uniquely identified based on increasing distance from the originating heavy atoms. These segments and identifiers are extended through bonds to generate larger substructures and corresponding identifiers. Finally, all of these substructures are hashed into a fixed length binary fingerprint (for this case study 1024 was used) which can be used for training. We employed SVR and Random Forest Regressor (RFR) algorithms within Scikit-learn python library for training our model.

%Graph Convolution (GC): The basic idea behind GC is that it first encodes the system as a graph with each of the atoms corresponding to a node and the bonds as edges. For each node and edge a feature vector is assigned consisting of 75 atomic features including atom type, degree, implicit valence, and formal charge. A neighbour list is also generated for each node. This is in contrast to ECFP as connectivity becomes an inherent feature of this technique. These vectors are then passed through a convolutional neural network consisting of one convolutional hidden layer and a fully connected layer. While more hidden layers are possible, in the interest of computational time a model with one layer was fully optimized.

%Weave: Like GC, Weave is an adaptive graph-based model that treats molecules as undirected graphs. Weave also uses the same features as GC for updating the atom feature vector. However, instead of doing convolutions locally (among neighbor lists), it includes longer range pair wise interactions through global convolutions. This is done by introducing a pair feature vector which encodes features such as bond properties, graph distance and ring information. Therefore, the complexity of the Weave model allows for the encoding of more structural information than the other two techniques, but at the expense of greater computation time. 

%------------------------------------------------------------------------
\subsection{Model Optimization}

The procedure for model optimization for each featurization algorithm was to separate a test set and optimize the model on training and validation sets using 5-fold cross-validation. Each hyperparameter was input as a range over which the model was initialized and trained in separate for loops and the RMSE performance over each iteration was analyzed. The hyperparameters yielding the best cross-validation score were taken to be those that produced the optimized model. Once all three models were optimized, they were trained on the full train-validation set and performance was compared on the separated test set (which was kept constant to allow for direct performance comparison). Performance was evaluated by calculating the RMSE, mean absolute percent error (MAPE), and visualized using parity plots. The equations for calculating the RMSE and MAPE are as follows:

\begin{equation}
    \text{RMSE} = \sqrt{\sum_i^N \frac{(\hat{y_i}-y_i)^2}{N}}
\end{equation}

\begin{equation}
    \text{MAPE} = \frac{100\%}{N} \sum^N_i \frac{y_i-\hat{y_i}}{y_i}
\end{equation}

% The optimization procedure for each featurization technique is discuand discuss their performance with on the separated test set. For both ECFP and GCNN only datapoints with boiling points below 700 K were considered. In the Weave evaluation section the results of using the full temperature range are discussed, and further motivation for introducing the 700 K cutoff are provided.


\begin{table}
\begin{center}
\small
\begin{tabular}{|c|c|c|}
\hline
Optimized Hyperparameters (architecture)& GC & Weave\\
\hline\hline
\# of Features & 75 & 75\\
\# of Hidden Layers & 1 & 2\\
Channel Width & 100 & -\\
Batch Size & 25 & 60\\
Learning Rate & 0.01 & 0.001\\
Epochs & 20 & 100\\
\hline
\end{tabular}
\end{center}
\caption{Comparison of the optimized parameters for GCNN and Weave.}
\label{table:hyp_para}
\end{table}

\subsection{Model Performance Evaluation}

In this section, we briefly discuss qualitative evaluation results for a ECFP featurization, graph convolution model, and weave model.

The ECFP algorithm computes hashed out fixed length matrices for each molecule which we fed into two different regressors: support vector regressor (SVR) and random forest regressor (RFR). Two parameters were optimized for the case of SVR: regularisation parameter (C) and width of no penalty ($\epsilon$). Similarly, two parameters were optimized for the case of RFR: depth of trees and number of estimators. However, even the best fit gives high RMSE and MAPE for both melting and boiling points. Both SVR and RFR models from Scikit-learn python library were relatively quick in training time, indicating simplicity and a deficiency of essential features to explain physical properties of molecules in ECFP featurizer. Moreover, ECFP misses out on the connectivity of the structure which could play an important role in determining the critical points of the system.

%\begin{marginfigure}
%   \includegraphics[width=\linewidth]{figures/ecfp_report.png}
%   \caption{Test set parity plot and performance of A) melting point and B) boiling point predictions using the ECFP featurizer.}
%\label{fig:ecfp}
%\end{marginfigure}

When optimizing the graph convolutional model, hyperparameters were individually tuned in the following order: channel width, batch size, learning rate, and number of epochs. For each of these parameters a range of potential values were iterated over, with the others held fixed. The optimal hyperparameter was selected considering both the 5-fold CV score and the computational weight of selecting the value. For example, if the score plateaus as a function of a given parameter, but increasing parameter means heavier computation time, then the lowest value where the score converges is chosen. In some cases the trends were hard to rationalize and the optimal value was not obvious, so some user judgement was necessary. %For example, the plot of channel width versus 5-fold CV score is shown in Fig. \ref{fig:opt}. A width of 100 was chosen from this plot as afterwards there appears to be some overall convergence in the error for both Melting and Boiling points. A summary of the obtained optimized parameters for this model are listed in Table \ref{table:hyp_para}.

%\begin{marginfigure}
%   \includegraphics[width=\linewidth]{figures/channel_width_opt.png}
%   \caption{Cross validation score as a function of channel width. It is observed that after 100 there appears to be some convergence in the error estimate}
%\label{fig:opt}
%\end{marginfigure}


Using these optimized parameters, GCNN was evaluated on the test set. RMSE scores of 51.42 and 52.20 for melting and boiling points respectively were obtained. For melting point most of the error stems from points on the right side of the figure where the model under predicts the melting points. Similarly, for boiling point most of the incorrectly predicted points lie to the top left of the figure indicating over prediction of the boiling points. This could be due to the inclusion of structure in the model overestimating the strength of the bond interactions. In comparison to ECFP, GCNN resulted in a much improved model.

%\begin{marginfigure}
%   \includegraphics[width=\linewidth]{figures/gcnn.png}
%   \caption{Test parity plot and performance of A) melting point and B) boiling point predictions for temperatures using the GCNN featurizer.}
%\label{fig:gcnn}
%\end{marginfigure}



For the weave model, melting point and boiling point predictions were significantly improved with RMSE errors of 50.31 and 48.06. The optimized hyperparameters on this limited dataset are given in Table \ref{table:hyp_para}. In comparison to the other models which only used temperatures below the 700 K threshold, Weave's metrics were marginally better than GCNN. Time taken to compute train model and prediction took 12 minutes which is higher than Graph convolution due to increased complexity of pair-wise interactions in updating atom features. The RMSE and MAPE are the least with respect to GCNN and ECFP featurizers. The source of improvement is most likely the pairwise features of Weave allowing a more accurate representation of the interatomic interactions than the other models.

%\begin{marginfigure}
%   \includegraphics[width=\linewidth]{figures/weave_report.png}
%   \caption{Test set parity plot and performance of A) melting point and B) boiling point predictions using the Weave featurizer.}
%\label{fig:weave}
%\end{marginfigure}

%------------------------------------------------------------------------
\subsection{Conclusions}
In this case study we have demonstrated the usage of ECFP, GCNN, and Weave models to simultaneously predict melting points and boiling points on a wide class of materials. We show that when using the truncated dataset ECFP demonstrated poor predictive capabilities, but GCNN and Weave showed promising performances. Weave was able to yield slightly lower RMSE than GCNN for both critical points, but this was at the expense of computation time during fitting. However, despite the increase in computational weight, we believe that Weave showed the best predictive power, and thus is the best model of the three featurization techniques. We attribute this improved performance to the inclusion of pair-wise interactions beyond nearest neighbours. Additionally, using a more systematic approach to optimizing the hyperparameters could open the door to even better displayed performances. %Continued use of Weave is a promising avenue for predicting melting and boiling points on a range of candidate electrolyte solvent materials with the purpose of rapid screening. Moreover, we believe that the framework presented here can be extended to train and fit models for other important electrolyte parameters in the ongoing effort of rapid identification of the most promising solvent substances. This case study demonstrates deep learning as an effective tool in predicting collective properties of materials and has the potential to have an impact on the overall speed of materials discovery, and accelerate the development of future green technologies.

\section{Liquid Crystals (Transition Temperature)}
Liquid crystals (LCs) are a state of matter which has properties between those of conventional liquids and those of solid crystals. One industrially relevant application of LCs is liquid crystal displays, which are a form of visual displays used in electronic devices wherein a layer of a liquid crystal is sandwiched between two transparent electrodes. A thermotropic nematic liquid crystal transitions to an isotropic liquid at the nematic–isotropic transition temperature (TNI), which depends on the molecular order of the mesophase. 
\end{document}